% ============================================================================
% 模型假设 —— 每条假设必须附「合理性依据」，光列不证等于没写。
%
% 【本节的高价值内容：口径决策表的披露版】
%   P1 口径审定（G1）产出的决策表，在这里给出面向读者的版本：
%   采用的口径 + 被否口径 + 一句话定量归谬（例：若按 B 口径，阈值将偏高两个
%   数量级，与题面某条件矛盾）。
%   这条是 loop1-r3 判负的直接教训：混合口径 = 方法级 fail；
%   而单一口径 + 反事实量化恰恰是最能加分的地方。
% ============================================================================
\section{模型假设}

\begin{enumerate}
  \item \textbf{假设一。} 依据：……
  \item \textbf{假设二。} 依据：……
\end{enumerate}

\subsection{口径决策与反事实核验}
\begin{table}[htbp]
  \centering
  \caption{关键口径决策表}
  \begin{tabularx}{\linewidth}{llX}
    \toprule
    决策点 & 采用口径 & 被否口径与定量归谬 \\
    \midrule
     &  &  \\
    \bottomrule
  \end{tabularx}
\end{table}
